\scanpage{10}\mspage{9}
\makeatletter\tagsleft@true\makeatother
of \uline{compatibility} in $\unsure{\mathcal{M}|M}$ is an important \deleted{notion} relation, which is \unsure{added|added} to the relations $\underset{\mathcal{M}}{\leq}$, $\underset{\mathcal{M}}{\ll}$. Morally, to recover the workshop (up to unique isomorphism) it suffices to know on $\mathcal{M}$ these two order relations, and the compatibility relation.
% slanted
\margin{at least in the tame case}

The \deleted{\unsure{following}} relations \inserted{(compatibility, disjointness)} are so important, that each deserves a terse notation. I propose
\[
\deleted{F\parallel G}\qquad F\parallel G \tag*{(1.14)}
\]
for disjointness, and
\[
\deleted{F\unsure{\mathrel{\tikz[baseline=0pt,x=1em,y=1em]{\draw (-.38,.62)--(.38,.25)--(-.43,-.10); \draw (.34,.62)--(-.32,.25)--(.38,-.10);}}|\mathrel{\tikz[baseline=0pt,x=1em,y=1em]{\draw (-.44,.26)--(0,.58)--(.44,.26)--(0,-.06)--cycle; \draw (0,.58)--(0,-.06);}}}G}\qquad F\mathrel{\tikz[baseline=0pt,x=1em,y=1em]{\draw (-.44,.26)--(0,.58)--(.44,.26)--(0,-.06)--cycle; \draw (0,.58)--(0,-.06);}} G\quad\text{or}\quad
\tikz[baseline=(fg.base)]{\node[inner sep=1pt] (fg) {$F\mathrel{\tikz[baseline=0pt,x=1em,y=1em]{\draw (-.38,.62)--(.38,.25)--(-.43,-.10); \draw (.34,.62)--(-.32,.25)--(.38,-.10);}} G$}; \draw ([shift={(70:2.12em and 1.40em)}]fg.center) arc[start angle=70,end angle=410,x radius=2.12em,y radius=1.40em];}
\tag*{(1.15)}
\]
for compatibility (NB $\mathrel{\tikz[baseline=0pt,x=1em,y=1em]{\draw (-.38,.62)--(.38,.25)--(-.43,-.10); \draw (.34,.62)--(-.32,.25)--(.38,-.10);}}$ is quicker \unsure{:|to} write than $\mathrel{\tikz[baseline=0pt,x=1em,y=1em]{\draw (-.44,.26)--(0,.58)--(.44,.26)--(0,-.06)--cycle; \draw (0,.58)--(0,-.06);}}$, $\unsure{\mathrel{\tikz[baseline=0pt,x=1em,y=1em]{\draw (-.12,.56)--(-.48,.25)--(-.12,-.02); \draw (.10,.56)--(.49,.25)--(.13,-.03); \draw (-.32,.51)--(-.59,.25)--(-.28,.18)--(.20,.30)--(.49,.25); \draw (-.28,.18)--(.02,.45); \draw (.26,.55)--(.50,.25);}}|\deleted{\mathrel{\tikz[baseline=0pt,x=1em,y=1em]{\draw (-.12,.56)--(-.48,.25)--(-.12,-.02); \draw (.10,.56)--(.49,.25)--(.13,-.03); \draw (-.32,.51)--(-.59,.25)--(-.28,.18)--(.20,.30)--(.49,.25); \draw (-.28,.18)--(.02,.45); \draw (.26,.55)--(.50,.25);}}}}$ but is less suggestive of the combinatorial idea expressed.

We have said nothing so far about the refinement relation $\deleted{\ll}\preccurlyeq$ on $\unsure{\mathcal{F}|\mathfrak{F}}$. It is formally of a more delicate nature, and is the most specific notion in the ``Algebra of figures''. The geometric meaning becomes clear through the \unsure{examples|example}, and the axioms that \unsure{these-|these}%
